\documentclass[11pt,a4paper]{article}
\usepackage[T1]{fontenc}
\usepackage[utf8]{inputenc}
\usepackage[main=italian,provide=*]{babel}
\usepackage{amsmath,amssymb}
\usepackage{geometry}
\geometry{margin=2.3cm,top=2.5cm,bottom=2.7cm}
\usepackage{booktabs}
\usepackage[table,dvipsnames]{xcolor}
\usepackage{tcolorbox}
\tcbuselibrary{skins,breakable}
\usepackage{titlesec}
\usepackage{enumitem}
\usepackage{needspace}
\usepackage{fancyhdr}
\usepackage{graphicx}
\usepackage{hyperref}
\hypersetup{colorlinks=true,linkcolor=Sepia,citecolor=Sepia,urlcolor=Sepia}
\newcommand{\Tr}{\operatorname{Tr}}
\newcommand{\ket}[1]{\lvert #1 \rangle}
\newcommand{\bra}[1]{\langle #1 \rvert}

\definecolor{inkblue}{HTML}{1B3A5C}
\definecolor{lightblue}{HTML}{EAF1F8}
\definecolor{accent}{HTML}{B0562A}
\definecolor{softgray}{HTML}{6B6B6B}
\definecolor{okgreen}{HTML}{2E6B3E}
\definecolor{okgreenbg}{HTML}{EAF5EC}

\titleformat{\section}
  {\normalfont\Large\bfseries\color{inkblue}}
  {\thesection}{0.6em}{}[{\color{inkblue!40}\titlerule}]
\titlespacing{\section}{0pt}{0.9em}{0.4em}
\titleformat{\subsection}
  {\normalfont\large\bfseries\color{accent}}
  {}{0em}{}
\titlespacing{\subsection}{0pt}{0.6em}{0.2em}

\pagestyle{fancy}
\fancyhf{}
\renewcommand{\headrulewidth}{0.4pt}
\fancyhead[L]{\small\color{softgray}Quantum simulation (Trotter) sotto rumore (anello)}
\fancyhead[R]{\small\color{softgray}\thepage}
\fancyfoot[C]{\small\color{softgray}Tesi triennale, Samuele --- Relatore: Prof.\ Alessandro Chiesa}

\newtcolorbox{keyeq}[1][]{
  colback=lightblue, colframe=inkblue!70, boxrule=0.6pt,
  arc=2pt, left=8pt, right=8pt, top=3pt, bottom=3pt, #1
}
\newtcolorbox{okbox}[1]{
  colback=okgreenbg, colframe=okgreen!60, boxrule=0.6pt,
  arc=2pt, left=7pt, right=7pt, top=2pt, bottom=2pt,
  fonttitle=\bfseries\color{okgreen}, title={#1}
}

\title{\vspace{-1.5em}\color{inkblue}Quantum simulation (Trotter) sotto rumore --- trimero ad anello\\[0.2em]
\Large\normalfont\color{softgray}Due scenari: punto VQE e punto $R_0$}
\author{Note di lavoro --- Tesi triennale, Samuele\\ Relatore: Prof.\ Alessandro Chiesa}
\date{}

\begin{document}
\sloppy
\maketitle
\thispagestyle{fancy}
\vspace{-2.0em}

File di codice:
\texttt{trotter\_\allowbreak rumoroso\_trimero\_anello.py} (mirror
strutturale di \texttt{trotter\_rumoroso\_dimero.py}, ma generico rispetto
al punto di lavoro e allo stato iniziale). Modello di rumore:
\texttt{noise\_model\_trimero\_anello.py}, valori di riferimento
\texttt{ibm\_torino}.

\tableofcontents

\section{Obiettivo e una scelta esplicita di scope}

A differenza del dimero, dove un unico punto di lavoro (``test 2'')
serve sia il VQE sia la dimostrazione Trotter, \textbf{Parte 1
dell'anello usa due punti diversi} per le due cose:

\begin{keyeq}[title={Due scenari, entrambi trattati qui}]
\begin{itemize}[leftmargin=1.3em,itemsep=0.2em]
\item[\textbf{A}] Punto VQE ($b=b_c=2.4,\,D=0.15$, mode ``B''): stato
preparato dall'ansatz $W$-2q.6 (Stadio 2), poi $N$ passi di Trotter.
Mirror esatto dello schema del dimero.
\item[\textbf{B}] Punto $R_0$ ($b=0.05,\,D=1.93$), \textbf{dichiarato
provvisorio} in Parte 1
(\texttt{quantum\_\allowbreak simulation\_\allowbreak trimero\_anello\_\allowbreak applicazione.tex}):
nessuna preparazione VQE, si parte da $\ket{000}$ --- stesso stato e
stesso punto della dimostrazione Trotter di Parte 1, qui reso rumoroso.
\end{itemize}
Scelta dell'utente: trattare \textbf{entrambi}, non scegliere fra i due.
\end{keyeq}

\section{Strategia di transpilazione a blocchi (scritta qui per la prima volta)}

Segnalato all'apertura della Parte 2 sull'anello: né
\texttt{circuito\_\allowbreak correlazioni\_\allowbreak trimero\_anello.py}
né \texttt{trotter\_trimero\_anello.py} contenevano alcun riferimento a
\texttt{BASIS\_GATES} --- verificato, non assunto. Stessa insidia già
affrontata sul dimero: transpilare l'\emph{intero} circuito collasserebbe
gli $N$ passi identici in un costo costante. Soluzione, identica nello
spirito: il singolo passo esterno (scambio $H_\text{ex}$, poi campo, poi
DM $H_\text{DM}$ --- stesso ordine di \texttt{trotter\_circuit} in Parte
1) viene transpilato \textbf{una sola volta} in base $\{rz,sx,x,cx\}$,
poi il blocco già compilato è composto $N$ volte.

\begin{okbox}{Corretto: livello di ottimizzazione, dimostrato minimo (non assunto)}
Richiesta esplicita del relatore (\texttt{domande\_relatore.md},
sez.~12): transpilare per \textbf{minimizzare} i gate, non usare il
default. Sul dimero questo fu dimostrato con un confronto esplicito
($10\to3$ CNOT, fattore $3.3\times$). Su questo documento,
inizialmente, \texttt{optimization\_level=3} era stato usato per
convenzione ereditata dal dimero, \textbf{senza} il confronto
esplicito --- gap trovato in una verifica indipendente successiva,
colmato qui:
\begin{center}
\renewcommand{\arraystretch}{1.25}
\begin{tabular}{@{}lcc@{}}
\toprule
\textbf{Livello} & \textbf{CNOT/passo (Trotter)} & \textbf{CNOT (ansatz VQE)}\\
\midrule
0 & 30 & 6\\
1 & 30 & 6\\
2 & 15 & 6\\
3 & 15 & 6\\
\bottomrule
\end{tabular}
\end{center}
\textbf{Confermato}: livello 3 (usato in tutta la pipeline) è
effettivamente minimo per il passo di Trotter (fattore $2\times$ su
livello 0/1). Sull'ansatz VQE nessun beneficio a nessun livello (6
CNOT costanti) --- risultato legittimo: la struttura a blocchi $W$
sequenziali su coppie di qubit diverse non offre ridondanza da fondere,
non un errore o un'occasione mancata.
\end{okbox}

\section{Risultati: conteggio dei gate}

\begin{center}
\renewcommand{\arraystretch}{1.3}
\begin{tabular}{@{}lccc@{}}
\toprule
\textbf{Scenario} & \textbf{CNOT preparazione} & \textbf{CNOT per passo} & \textbf{CNOT totali}\\
\midrule
A (VQE) & 6 & 15 & $6+15N$\\
B ($R_0$) & 0 & 15 & $15N$\\
\bottomrule
\end{tabular}
\end{center}

\begin{okbox}{Verificato}
Crescita lineare confermata per $N\in\{1,2,5,10,20\}$ in entrambi gli
scenari, coincidenza esatta con la formula attesa --- nessuna fusione
silenziosa dei passi.
\end{okbox}

\section{Cross-check indipendente (rumore nullo)}

Stessa struttura del dimero: fedeltà dal circuito Qiskit confrontata con
un calcolo classico via potenza di matrice del singolo passo (numpy
puro), cammino di calcolo indipendente dalla transpilazione a blocchi.

\begin{center}
\renewcommand{\arraystretch}{1.25}
\begin{tabular}{@{}lccc@{}}
\toprule
\textbf{Scenario} & $\boldsymbol{N}$ & $\boldsymbol{F}$ \textbf{(numpy)} & $|$\textbf{diff}$|$ \textbf{vs Qiskit}\\
\midrule
A (VQE) & 5 & $0.9844752881$ & $4.4\times10^{-9}$\\
A (VQE) & 80 & $0.9999521509$ & $2.6\times10^{-10}$\\
B ($R_0$) & 5 & $0.3567981074$ & $3.1\times10^{-15}$\\
B ($R_0$) & 80 & $0.9960968282$ & $1.6\times10^{-13}$\\
\bottomrule
\end{tabular}
\end{center}

\begin{okbox}{Verificato}
Entrambi gli scenari convergono a $F\to1$ per $N$ crescente a rumore
nullo (coerente con la convergenza già documentata in Parte 1). Lo
scenario B converge più lentamente ($F=0.996$ a $N=80$, contro
$F=0.99995$ dello scenario A): coerente con Parte 1, dove $D/J\simeq1.9$
al punto $R_0$ (non piccolo) richiede $N\gtrsim300$ per infedeltà
$<1\%$, contro il regime più mite del punto VQE.
\end{okbox}

\section{Il compromesso Trotter/rumore: localizzazione di \texorpdfstring{$N^*$}{N*}}

\begin{figure}[h]
\centering
\includegraphics[width=0.98\linewidth]{figure/fig_trotter_rumoroso_trimero_anello.pdf}
\caption{\textit{Fedeltà rumorosa $F(N)$ al rumore di riferimento
(\texttt{ibm\_torino}), per entrambi gli scenari. In entrambi i casi il
massimo è a un $N$ finito --- il compromesso fra errore di Trotter
(decresce con $N$) e rumore accumulato (cresce con $N$, $15$ CNOT a
passo) si manifesta concretamente, non solo previsto teoricamente.}}
\end{figure}

\begin{okbox}{Corretto (griglia $N$ infittita in Stadio 5)}
Valore originale di questo documento: $N^*_A\approx2$,
$N^*_B\approx10$, da una griglia $N\in\{1,2,5,10,15,20,30,40,60,80\}$
che \textbf{saltava} $N=3$ e $N=4$. Rieseguendo con una griglia continua
$N=1,\ldots,20$ (Stadio 5, \texttt{scan\_parametri\_rumore\_trimero\_
anello.py}) sono emersi i veri massimi: $N^*_A=3$
($F(N^*)=0.700335$, contro $0.697$ al vecchio $N=2$ --- vicini ma non
uguali) e $N^*_B=9$ ($F(N^*)=0.381896$, contro $0.380$ al vecchio
$N=10$ --- praticamente indistinguibili). La conclusione qualitativa
(due $N^*$ diversi, entrambi finiti) non cambia; i valori numerici sì,
e sono stati corretti qui e nella figura.
\end{okbox}

\begin{okbox}{Verificato}
$N^*_A=3$ ($F(N^*)=0.700335$), $N^*_B=9$ ($F(N^*)=0.381896$) ---
\textbf{diversi fra loro}, non assunti identici al dimero né fra i due
scenari stessi. Lo scenario B mostra inoltre un comportamento non
monotono a $N$ piccolo ($F$ crolla a $N=2$ prima di risalire): a rumore
nullo, $F(N{=}2)=0.0017$ è già piccolissima per il solo errore di
Trotter (regime non perturbativo a $D/J\simeq1.9$), quindi il crollo a
$N=2$ sotto rumore riflette innanzitutto un errore di Trotter grande a
quel passo, non un effetto del rumore in sé.
\end{okbox}

\section{Cosa non è ancora stato controllato qui}

Nessuno scan sistematico sui parametri di rumore (Stadio 5, non ancora
aperto): qui $N^*$ è localizzato solo al punto di calibrazione di
riferimento. Nessuna correlazione dinamica (Stadio 4): qui l'osservabile
è la fedeltà rispetto allo stato fisico evoluto, non un correlatore a
due tempi.

\begin{okbox}{Punto lasciato esplicitamente aperto, non verificato}
Lo Stadio 2 ipotizzava che la degradazione di preparazione
($1-F\sim3.0\%$) entrasse nella fedeltà totale di Trotter come
\textbf{contributo additivo e costante in $N$}. Questo documento
\textbf{non verifica} tale decomposizione: calcola $F(N)$ nel suo
complesso, non i singoli contributi (preparazione, errore di Trotter,
rumore accumulato) separatamente. Precedente rilevante sul dimero
(\texttt{log\_decisioni.md}, sessione 5 settembre): la stessa
decomposizione additiva era stata verificata lì per un \emph{correlatore}
(Stadio 4), poi estesa per analogia alla fedeltà di Trotter (Stadio 3)
\textbf{senza} una derivazione indipendente per quella metrica --- un
errore di ragionamento poi corretto esplicitamente. Per non ripeterlo
qui: la frase dello Stadio 2 va quindi letta come un'ipotesi di lavoro
non ancora messa alla prova, non come un fatto già stabilito.
\end{okbox}

\section{Considerazioni e cosa serve al passo successivo}

Lo Stadio 4 (correlatori dinamici sotto rumore) userà lo stesso blocco
di passo esterno transpilato qui (\texttt{build\_step\_block}), inserito
nel circuito Hadamard test con qubit ancilla --- il costo per passo (15
CNOT) resta lo stesso, a cui si aggiungerà il costo del circuito di
misura.

\begin{okbox}{Correzione (aggiunta dopo la chiusura dello Stadio 4)}
Questa sezione prevedeva, quando scritta, che allo Stadio 4 andassero
verificati entrambi gli scenari A e B. \textbf{Previsione non
confermata}: Parte 1 ha esplorato i correlatori dinamici \emph{solo} al
punto VQE (scenario A) --- il punto $R_0$ non è mai stato usato per i
correlatori, nemmeno in Parte 1. Lo Stadio 4 verifica quindi solo lo
scenario A, correttamente, non per un'omissione ma perché lo scenario B
non ha un analogo da riprodurre per questa osservabile. Vedi
\texttt{risultati\_\allowbreak correlatori\_\allowbreak rumorosi\_trimero\_anello.tex},
Sez.\ ``Cosa non è ancora stato controllato qui''.
\end{okbox}

\begin{okbox}{Sintesi in una riga}
La strategia di transpilazione a blocchi, scritta qui per la prima volta
per l'anello, funziona identicamente in entrambi gli scenari (crescita
lineare confermata); il compromesso Trotter/rumore produce un $N^*$
finito ma \textbf{diverso} nei due scenari ($N^*_A=3$, $N^*_B=9$),
segno che il punto di lavoro conta quanto il livello di rumore.
\end{okbox}

\needspace{4\baselineskip}
\section*{File prodotti}
\addcontentsline{toc}{section}{File prodotti}

\texttt{trotter\_\allowbreak rumoroso\_trimero\_anello.py},
\texttt{validate\_\allowbreak trotter\_\allowbreak rumoroso\_trimero\_anello.py},
\texttt{generate\_figure\_\allowbreak trotter\_\allowbreak rumoroso\_trimero\_anello.py},
questo documento.

\end{document}
